\chapter{Introduction: On the measure of intelligence}
\pagenumbering{arabic}

\begin{flushleft}
    \textit{On the measure of intelligence}\cite{cholletMeasureIntelligence2019} is a 2019 paper that aims to give an actionable definition of intelligence and proposing a new benchmark for AI. We will discuss the purpose and implications of this definition and how the benchmark has evolved since the inception of this paper.
\end{flushleft}

\begin{flushleft}
    First of all, we intuitively understand what intelligence is, when someone points to a person and calls them intelligent, we know what they mean by that. But arriving at a conclusive all-encompassing definition of what exactly entails intelligence is incredibly difficult, and we will not reach a consensus on that. However, note that it was explicitly said to be an \textit{actionable} definition, meaning that the purpose of this definition is to be useful towards measuring intelligence. In order to provide context, we'll have to delve into a bit of psychology.
\end{flushleft}


\section{Describing the mind}

\subsection{Early interpretations of the mind}

    \begin{flushleft}
        Since the inception of psychology, there have been two opposing views that describe the human mind\cite{spelkeCoreKnowledge2007}: a relatively static assembly of modules developed through evolution and a blanks slate that turns experience into knowledge and skills. 
    \end{flushleft}

    \begin{flushleft}
        The first point of view, rooted in evolutionary psychology, strongly influenced the work of early AI researchers, which viewed intelligence as static programs, relying on logical operators and database-like memory. Even know, we see AI being reduced to task performers. 
    \end{flushleft}

    \begin{flushleft}
        Contrasting with evolutionary psychology, we have the Locke's \textit{Tabula Rasa}, a vision of the mind as a blank slate that turns experience into skills. One might come up with the metaphor of the random initialization of a Neural Network as the blank slate which learns from training data. However, we would be ignoring the surrounding context, that being the work put by the engineers carefully designing the architecture and metrics to tackle the task, esentially enclosing the domain. Again, we would be falling in the same pitfalls as those early AI researchers.
    \end{flushleft}

    \begin{flushleft}
        In reality, developmental psychology show us that none of these views precisely describe the mind. We do have some "hard-coded programs", so called cognitive priors that will be thoroughly discussed later, but we still acquire most skills and knowledge through our lifetime. And yet, we are not a blank slate capable of learning anything through experience. Ultimately, what makes us intelligent are not those skills nor knowledge, but the process by which we acquire them and are able to adapt to unseen situations.
    \end{flushleft}

    \subsection{Separating the innate from the acquired: insights from developmental psychology}
    \subsection{The spectrum of generalization: robustness, flexibility, generality}
    \subsection{The meaning of generality: grounding the g factor}

    

\section{Evaluating intelligence}


    \begin{flushleft}
        Our next point is how to evaluate intelligence, but before discussing how to actually do that, we'll first discuss why it is important in the first place to that our AI system is intelligent. Then, we'll take a look at how we evaluate an AI's performance at completing well-defined tasks and going deeper into what intelligence entails.
    \end{flushleft}

    \subsection{Why it's important to do so}

    \subsection{Narrow AI evaluation}
        \begin{flushleft}
            Historically, these have been some of the most succesful approaches at quantifying the performance of an AI system:
        \end{flushleft}

        \begin{itemize}
            \item Human review: having actual humans judging the system's response and scoring it, though it is rarely used in practice since it is quite expensive.
            \item White-box analysis: inspecting the implementation of the system to determine its input-output response and score it.
            \item Peer confrontation: having the system compete against either other AIs or humans.
            \item Benchmarks: having the system produce outputs for a “test set” of inputs (or environments) for which the desired outcome is known, and score the response.
        \end{itemize}

        \begin{flushleft}
            Benchmarks have demonstrated to be particularly useful due to their characteristics: setting clear goals and adopting objective measure of performances that are shared across the reasearch community. However, heavily optimizing for some metrics has tradeoffs. For example, winners of Kaggle competitions are too specialized to be useful in real world environments.
        \end{flushleft}

        \begin{flushleft}
            However, as we discussed in the previous section (Describing the mind), intelligence can't simply be peforming well at completing some tasks, and is closer to the process of skill acquisition. 
        \end{flushleft}
    
    \subsection{AI effect}
    \subsection{Measuring broad abilities and general intelligence: the psychometrics perspective}
    \subsection{Measuring the right thing: evaluating skill alone does not move us forward}
    cognitive priors

\section{Defining intelligence: a formal synthesis}
\section{Discussion}

Follow ups on modern LLMs (foundational and agentic models)?
